\subsection{محدودیت‌ها و تحلیل شکست}
با وجود انعطاف‌پذیری روش، محدودیت‌های فنی زیر قابل تامل است:
\begin{itemize}
    \item \textbf{تخریب مسیر انتشار گرادیان با حذف لایه‌ها در مدل‌های با عمق کم:} در مدل‌های فشرده نظیر \lr{DistilBERT} که دارای تنها ۶ لایه هستند، حذف حتی یک لایه منجر به از دست رفتن بیش از ۱۶ درصد ظرفیت پردازشی مدل شده و می‌تواند رانش بازنمایی جبران‌ناپذیری ایجاد کند.
    \item \textbf{انباشت خطای صفر کردن کانال‌ها در شبکه‌های دیکودر:} در مدل‌های علّی نظیر \lr{GPT-2}، صفر شدن کانال‌های ورودی لایه توجه بدون تغییر متناظر در لایه خروجی می‌تواند تعادل وزن‌های پیش‌بینی توکن‌های بعدی را بر هم زند و افت نسبی دقت را به همراه داشته باشد.
    \item \textbf{نیاز به کالیبراسیون داده‌محور جهت انتخاب بهینه لایه‌ها:} انتخاب لایه‌ها برای حذف صرفاً بر اساس الگوی تصادفی میانی نیازمند ارزیابی دقیق حساسیت باقیمانده بر روی دادگان ارزیابی است تا از حذف تصادفی لایه‌های حاوی اطلاعات ساختاری زبانی پیشگیری شود.
\end{itemize}
